\chapteropen
\chapter{A history of multi-wavelength observations of galaxy clusters}
\label{app:mw_obs}

X-ray emissions from the hot gas in the \gls{icm} in clusters were first definitively detected in the 1970s; \eg \citet{Gursky_1971} details the first X-ray detection in the Coma Cluster by the Uhuru satellite. Detections have steadily improved over recent decades. For example, in the 1990s, the \gls{ROSAT} produced the first detailed surface brightness maps of clusters like Perseus, Coma and Virgo \citep{Boehringer_1993,Bohringer_1994}. This was followed by Chandra and XMM-Newton in the 2000s, which revealed \glspl{cold-front} \citep{Markevitch_2000}, \glspl{AGN-driven-cavity} \citep{Fabian_2000}, and \glspl{shock-front} \citep{Vikhlinin_2001} in the \gls{icm}. More recently, the short-lived \gls{jaxa} Hitomi mission measured turbulence in the Perseus cluster core \citep{HitomiCollaboration_2016}, and the ongoing \gls{eROSITA} mission is mapping thermodynamic properties of the \gls{icm} to redshifts beyond $z \sim 1$ \citep{Bulbul_2024}. 

The \gls{icm} is also visible at radio wavelengths due to \glslink{synchrotron-rad}{synchrotron emission} from relativistic particle populations, and so radio surveys by facilities such as \gls{lofar} \citep{Cuciti_2022} and the \gls{mgcls} have uncovered diffuse \glspl{radio-halo} and \glspl{radio-relic} tracing turbulent and shocked \gls{icm} regions \citep{Knowles_2022}. Most recently, the \gls{alma} in Chile, has revealed \gls{ram-pressure-stripping} effects on cluster galaxies \citep[\eg][]{Jachym_2014,Jachym_2017}. For example, observations of cold molecular gas tails in ``jellyfish'' galaxies \citep[\eg][]{Moretti_2020} in clusters (e.g., Coma, Virgo, A2744), have confirmed that \gls{icm}--galaxy interactions drive \gls{quenching} and morphology transformations.
